\documentclass[lettersize,journal]{IEEEtran}
\usepackage{amsmath,amsfonts}
\usepackage{amsthm}
\usepackage{algorithmic}
\usepackage{algorithm}
\usepackage{array}
\usepackage[caption=false,font=normalsize,labelfont=sf,textfont=sf]{subfig}
\usepackage{textcomp}
\usepackage{stfloats}
\usepackage{url}
\usepackage{verbatim}
\usepackage{graphicx}
\usepackage{cite}
\usepackage{multirow}
\hyphenation{op-tical net-works semi-conduc-tor IEEE-Xplore}
% updated with editorial comments 8/9/2021

\begin{document}

\title{Predicting Number of Clusters by Using Past Knowledge Transfer of Pre-trained Cluster Patterns
without Statistical Distance Analysis}

\author{Chidchanok Lursinsap ~\IEEEmembership{}
        % <-this % stops a space
\thanks{Department of Mathematics and Computer Science, Faculty of Science, Chulalongkorn University, 
Bangkok 10330, Thailand}% <-this % stops a space
\thanks{}}

% The paper headers
\markboth{Journal of \LaTeX\ Class Files,~Vol.~14, No.~8, August~2021}%
{Shell \MakeLowercase{\textit{et al.}}: A Sample Article Using IEEEtran.cls for IEEE Journals}

\IEEEpubid{0000--0000/00\$00.00~\copyright~2021 IEEE}
% Remember, if you use this you must call \IEEEpubidadjcol in the second
% column for its text to clear the IEEEpubid mark.

\maketitle

\begin{abstract}
Knowing the exact number of clusters and their centers in any dimension is the most important step to 
achieve the highest accuracy of any recognition problems. There are several efficient clustering methods 
proposed in the past to identify the clusters. These methods must specify the number of 
clusters (k-value) prior to the clustering process which deploys some statistical distance analysis with a
predefined decision threshold value. In this paper, a new approach based on the cognitive process of
the brain using the past knowledge and experience to brain to solve problems is proposed. The problem of 
identifying number of clusters is transformed into a problem of classifying data distribution patterns
with number of clusters as a target. The preliminary result achieved 98\% prediction accuracy. The 
theoretical upper bound in the number of pre-trained patterns is also given. An experiment on a real 
data set was also conducted with the discussion about outcome. 

\end{abstract}

\begin{IEEEkeywords}
Article submission, IEEE, IEEEtran, journal, \LaTeX, paper, template, typesetting.
\end{IEEEkeywords}

\section{Introduction}
\IEEEPARstart{C}{lustering} is a process to identify the number of groups or clusters of vector points 
and their locations in the space with any predefined targets and also the center of each group. 
The result of clustering can help improve 
the result of classification which is core of machine learning. If the number of clusters also its 
location are known a priori, then the minimum number of separating hyperplanes or radial basis functions 
can be setup without any trial-and-error process. To locate each cluster, most algorithms requires
a predefined parameters such as number of clusters, neighborhood size, branching factors before starting 
the clustering process. Table \ref{cluster-method-parameter} summarizes the parameters needed by some 
existing popular clustering methods implemented in package {\em scikit-learn}.  Therefore, knowing the 
actual values of these parameters is the most important step. But it is rather difficult to estimate the 
values of these parameters if it is more than 3 dimensions. Many proposed clustering algorithms 
employed the analysis of distances among data points with various metrics to determine the gaps among
potential clusters. Furthermore, the clustering process can be further improved by applying the concept of 
hierarchical grouping to merge some nearest clusters so that some false negative and redundant clusters 
can be eliminated. Although the approaches of distance analysis with deterministic metrics can produce a 
correct number of clusters in most cases, these approaches are ad hoc to a specific circumstance
and the accuracy depends strongly on a predefined threshold value,
evaluation metrics, and types of analyzing distance. It is noticeable that the knowledge and experiences
in predicting number of clusters obtained from previously analyzed cases cannot be transferred to new 
cases to help achieve the high accuracy of the new cases.
\begin{table}
\label{cluster-method-parameter}
\caption{Summary of some popular clustering methods and required parameters.}
\begin{tabular}{|l|l|} \hline
\multicolumn{1}{|c|}{Methods} & \multicolumn{1}{c|}{Parameters} \\ \hline
K-means & Number of clusters \\
Affinity propagation & Damping, Sample preference \\
Mean shift & Bandwidth \\
Spectral clustering & Number of clusters \\
Ward hierarchical clustering & Number of cluster or Distance threshold \\
Agglomerative clustering & Number of clusters or Distance threshold \\
DBSCAN & Neighborhood size \\
HDBSCAN & Minimum number cluster membership \\
OPTICS & Minimum number cluster membership \\
Gaussian mixtures & Number of Gaussian functions \\
BIRCH & Branching factor, Threshold \\
Bisecting K-means & Number of clusters \\ \hline
\end{tabular}
\end{table}

\IEEEpubidadjcol
In this study, a new approach is suggested so that the previous knowledge and experiences in predicting
the number of clusters prior to finding the location of each cluster can be transferred to cope with
new circumstances. This approach is based on the concept of {\em transfer of Knowledge
in cognitive systems}. By applying this concept, there is no need to analyze the relations among the 
data points in terms of distances and no evaluation metrics are required to determine the clusters.
The more knowledge of data distribution patterns and number of existing clusters are encountered
and experienced, the more accurate result of prediction can be achieved. Dagmar \cite{Dagmar} discussed
how the brain uses its past experiences to guide future actions, including predictions. He also
summarized the strong emerging trend towards
the study on memory-based prediction in cognitive neuroscience. Hawkins and Blakeslee 
\cite{Hawkins} stated that brain stores experiences as well as remembers events with their relations 
it encountered and uses these pieces of information and knowledge to make prediction. This is the
basis of brain intelligence, perception, creativity, and consciousness. Bar \cite{Bar} studied how 
the brain uses a priori learned knowledge to generate predictions. He also summarized that the brain is
proactive which means that it continuously produces predictions from the stored knowledge and also
keeps extracting and storing the information from the inputs.

\IEEEpubidadjcol
To avoid the analysis of distances among data points, a set of cluster distribution patterns with
their corresponding number of clusters as their targets are
transformed into a set set of vector patterns and targets. These patterns with the targets are learned 
by a supervised neural network. The network after being trained is used to predict the number of clusters.
The detail of this approach will be discussed in Section \ref{concept}. This approach does not
establish any constrained on the limitation of number of clusters, number of data points in each clusters,
and number of dimensions of data points. However, if the leaning methods proposed in 
{\em discard-after-learn} \cite{Prem}\cite{Mongkhon} are used, the training of new knowledge for the 
network can be continuously added without losing the previously learned knowledge.

The paper is organized into 5 sections.


\section{Related Works and Concepts of Predicting Number of clusters}

This section summarized some interesting works relating to predicting the number of clusters based on 
distance analysis among data points.

\IEEEpubidadjcol

Hong Yu, Zhanguo Liu, Guoyin Wang \cite{Hong} applied the concepts of rough set and hierarchical clustering 
to find the number of clusters in two algorithms, DTRS (decision-theoretic rough set) and FACA-DTRS
(fast automatically clustering algorithm using
decision-theoretic rough set). They also proposed some metrics to evaluate the clusters in terms
of probability. However, a predefined threshold to make decision must be set up.
Wilson Calmon and Mariana Albi \cite{Calmon} proposed two methods based on classical Plackett–Luce model.
The number of clusters must be defined in advance before the clustering process. Probability mass function 
was applied to rank vectors.
Md Abdul Masud, Joshua Zhexue Huang, Chenghao Wei, Jikui Wang, Imran Khan, Ming Zhong \cite{Abdul}
transformed the problem of predicting the number of clusters into the problem of finding number of hills
in data distribution surface. The surface is constructed in the form of distance distribution captured
Gamma mixture models (GMMs). Akaike information criterion variant (AICc) was used to select the hills.
Jiye Liang, Xingwang Zhao, Deyu Li, Fuyuan Cao, Chuangyin Dang \cite{Liang} used Re´nyi entropy and 
complement entropy as a metric to determine number of clusters for numerical and categorical data. 
They considered the entropy within cluster and between clusters. In addition, a dissimilarity measure 
also introduced to the k-prototypes algorithm.
Avisek Gupta, Shounak Datta, Swagatam Das \cite{Gupta} introduced the concepts of Last Leap (LL) and 
the Last Major Leap (LML) to find the number of clusters. Both concepts are based on the analysis of 
distances among data points. The maximum number of clusters must be given before the analysis. 
Ilias Gialampoukidis, Stefanos Vrochidis, Ioannis Kompatsiaris, Ioannis Antoniou \cite{Gialampoukidis}
modified the classical DBSCAN algorithm by adding probability martingale into consideration. The new 
algorithm is called DBSCAN-Martingale. The probability is computed from the distance between data points. 
Seong S. Chaea, Janice L. DuBienb, William D. Warde \cite{Chaea} used several agglomerative clustering 
algorithms and various metric indexes to determine number of clusters. They suggested a method to
decide whether the index values generated by different algorithms agreed or disagreed.
Mrinal Kanti Bhowmik, Tathagata Debnath, Debotosh Bhattacharjee, Paramartha Dutta \cite{Bhowmik}
applied Coulomb's law of electrostatics to imitate electrostatic force to cluster pixels of 
an image for image segmentation. Each pixel is treated as a point with some electrostatic charge which 
can create an attraction force. E. Lippiello, S. Baccari, P. Bountzis \cite{Lippiello} proposed a method
called Susceptibility of the Similarity Matrix. The matrix is transformed into another triangular matrix. 
The degree of similarity of two data points is measured by Euclidean distance with some additional 
parameters. Pingxin Wang, Yiyu Yao  \cite{Wang} applied the method of erosion and dilation in morphology
to determine the clusters. This approach must use a {\em hard clustering method} which referred to
existing clustering methods to find all tentative clusters first. Then they applied morphology to merged 
clusters. Ramazan Unlu, Petros Xanthopoulos \cite{Ramazen} used consensus clustering algorithm and 
different metrics which are Silhouette (SH), Calinski–Harabasz (CH), Davies–Bouldin (DB), and
Consensus (CI) indices to estimate the correct number of clusters. This approach is simple. Each 
existing clustering algorithm is applied and the result is evaluated by the metrics. Any result produced
the highest metric value is interpreted as the correct one.
Ahmed Fahim \cite{Fahim} added a pre-processing step to estimate an initial tentative number of 
clusters by applying a density-based method which does not require a predefined number of clusters. 
This estimated number of clusters is used as a $k$ value for further classical $k$-mean clustering methods.
The classical DBSCAN (Density-based spatial clustering of application with noise) is used in
the pre-processing step. Wei Luo \cite{Luo} proposed an interesting approach different from other 
clustering methods. They perturbed the given set clusters by augmenting another set of clusters as 
working clusters whose the number of perturbing clusters must be predefined. If the working clusters are 
coincident with the given clusters, then the number of clusters should be equal to the number of working. 
Otherwise the actual number of given clusters may be incorrectly estimated.  
Zhenzhen He, Zongpu Jia,and  Xiaohong Zhang \cite{He} considered the problem of estimate the number of 
clusters with complex distribution and scattered patterns. They used Gaussian kernel density function to 
estimate maximum number of clusters and their centers. The number of clusters is determined from
the change of minimum distance between center points. 
Toan Nguyen Mau, YAasushi Inoguchi, and Van-Nam Huynh \cite{Toan} used dissimilarity measure between two
data points and Locality-Sensitive Hashing (LSH) to produce fuzzy initial cluster prediction. Then the 
method of iterative fuzzy clustering was used next to obtain the final number of clusters. 
Xiaolong Wang, Yiping Jiao, Shumin Fei \cite{Xiaolong} used density and watershed method to partition
data space into multiple regions. The number of regional centers is set as the initial number of clusters. 
Md Moshiur Rahman, Abdul Masudy, Badhan Mazumderz \cite{Rahman} initially partitioned the selected 
features of data to explore the compactness of clusters. The depth value of data points in each clusters 
in the initial partition is estimated by {\em simplicial depth method}. The number of clusters having maximum depth value is considered as the estimated number of clusters. 
Dong Sik Kim \cite{Dong} studied the effect of ratio of estimated number of clusters and size of data 
points on the derivation of actual number of clusters.
This ratio is used in the proposed algorithm to adjust the value of $k$ in classical
$k$-mean clustering method.
 
\section{Problem Formulation and Proposed Concept}
The objective of this paper is to identify the actual number of clusters from a given set of data
points in any $n$-dimensional space without deploying any density analysis, information analysis,
and distance analysis. Each data point is represented in the form of coordinates. The following 
constraints are imposed on the studied problems.
 problem
has the following constraints.
\begin{enumerate}
\item The number of data points in each given sample must be in between the minimum and maximum 
number of data points.
\item The number of dimensions is at least 1.
\item The size of each dimension of the space containing data points in a sample is limited by 
the maximum value of coordinate in the corresponding dimension.
\end{enumerate}
Since no analysis is involved in this problem, the following two main problems are investigated.
\begin{enumerate}
\item Since the number of data points in each sample is different, how to transform a set of 
$n$-dimensional data points in each sample into a vector of fixed 
dimensions, regardless to the number of data points in the sample.
\item How to identify the actual number of clusters from the set of vectors obtained from problem 1
without analyzing the statistical information of these vectors? 
\end{enumerate}
 
The concept to solve the problems comes from two directions. The first direction is based of the 
uniqueness of data distribution pattern forming clusters. The second direction is from the studied 
result in cognitive science in terms of applying past knowledge and experience of the brain to solve new 
problems as mentioned by Dagmar \cite{Dagmar}.

Solving the problem of identifying the actual number of clusters from a given set of $n$-dimensional
data points without any statistical analysis is based on the observation that the distribution pattern
of data points of each sample is obviously unique. This results in a unique set of clusters in the 
data  pattern. Thus, a many-to-one mapping function can be developed to directly
compute the number of clusters from each data distribution pattern. This function is possibly 
realized by a multi-layer supervised learning neural network. It can be trained by a set of samples
whose actual number of clusters is known prior to use it to identify the number of actual clusters of
any unseen data distribution pattern. These training samples are considered as the past stored 
knowledge and experience.

The overview of proposed concept is illustrated in Figure \ref{overview}. There are two main processes.
The first process is to pre-train a network by a set of synthetic data distribution patterns with the
number of actual clusters as the targets. The second step is to identify the number of clusters of a test 
pattern by using the pre-trained network. There are 6 main steps in the pre-training process as follows.

\vspace{0.2in}
\noindent
{\bf Pre-training Process}
\begin{enumerate}
\item Synthesizing pre-training samples having different number of clusters.
\item Making the number of data points in each sample equal by oversampling the data points to add extra
data points.
\item Representing data distribution pattern in each sample by a vector of fixed dimensions.
\item Reducing the dimensions of representative vector by using self-encoding network. This step
is to speed up the learning process. The encoder is obtained from this step.
\item Pre-training encoded synthetic samples by a 3-layered network.
\end{enumerate}

\vspace{0.2in}
\noindent
{\bf Test Process}
\begin{enumerate}
\item Making the number of data points in the test pattern equal to number of data points in training
pattern by oversampling the data points to add extra data points.
\item Representing data distribution pattern in each sample by a vector of fixed dimensions.
\item Use encode from the first process to encode the vector from step 2 of the second process.
\item Use Pre-trained network from the first process to identify the number of clusters.
\end{enumerate}

\begin{figure*}[t]
\centering
\includegraphics[scale=0.6]{overview.eps}
\vspace{-0.6in}
\caption{Overview of proposed concept to identify actual number of clusters based on
the observation of uniqueness of data distribution pattern and the finding of how the brain uses
the past knowledge and experience to solve new problem.}  
\label{overview}
\end{figure*}

\section{Details of Past-Knowledge-Transfer Clustering Algorithms}

The following notations are used in the proposed algorithms. Algorithm 1 is the overview of the proposed
method. The other algorithms explain the detail of each step.
\begin{enumerate}
\item A set of pre-learning samples, ${\bf S} = \{{\bf s}_{1}, {\bf s}_{2}, \ldots, {\bf s}_{n}\}$.
\item Maximum number of points, $M$, and minimum number of points, $N$, such that $N \leq |{\bf s}_{i}| \leq M$.
\item Number of dimensions, $d$, of each point in the sample set.
\item A set of points in each ${\bf s}_{i}$, ${\bf P}^{(i)} = \{{\bf p}^{(i)}_{1}, {\bf p}^{(i)}_{2}, 
\ldots, {\bf p}^{(i)}_{m}\}$; ${\bf p}^{(i)}_{j} \in R^{d}$, $N \leq m \leq M$. 
\item Maximum number of clusters, $C \geq 1$.
\item Predicted bit string ${\bf B}^{(i)} = [b^{(i)}_{1}, b^{(i)}_{2}, \ldots, b^{(i)}_{C}]$; 
$b^{(i)}_{j} \in$ $\{0, 1\}$ and $b^{(i)}_{j} = 0.5$ if there are $j$ clusters in sample ${\bf s}_{i}$.
\end{enumerate}
Note that the value of target bit $b^{(i)}_{j}$ is set to 0.5 in order to speed up the convergence of outcome.
This value 0.5 can be interpreted as a threshold value as well.

\vspace{0.2in}
\noindent
{\bf Algorithm 1: Overview of proposed pre-learning sample method}
\label{alg1}        
\begin{tabbing}
1. \= Synthesize a set ${\bf S}$ of $n$ samples whose each ${\bf s}_{i}$ has various \\
   \> number of clusters not more than $C$ and different point \\
   \> distributions. For each sample ${\bf s}_{i}$, set the pattern of bit \\
   \> string ${\bf B}^{(i)}$ corresponding to the number of generated \\
   \> clusters. (see Algorithm 2) \\
2. \> Add extra data points to each sample ${\bf s}_{i}$ to make every \\
   \> $|{\bf s}_{i}| = M$. (see Algorithm 3) \\
3. \> For each ${\bf s}_{i}$, arrange all coordinates of each ${\bf p}^{(i)}_{j} \in {\bf P}^{(i)}$ \\
   \> to form a single sample vectors, ${\bf v}_{i} \in R^{d \times M}$ representing \\
   \> sample ${\bf s}_{i}$. (see Algorithm 3) \\
4. \> For each ${\bf s}_{i}$, set a proper target bit value of each $b^{(i)}_{j}$. \\
5. \> Get one testing sample and apply Algorithms 3 to add \\
   \> extra data points and to generate a new test sample vector ${\bf t}$. \\ 
6. \> Reduce the dimensions of all samples ${\bf v}_{i}$ and the testing \\
   \> sample ${\bf t}$ to create a new a set of vectors ${\bf v}^{\prime}_{i}$ and ${\bf t}^{\prime}$ by \\
   \> using the same self-encoder network. (see Algorithm 4) \\
7. \> Learn sample ${\bf s}_{i} \in {\bf S}$ with their corresponding target bit \\
   \> string ${\bf B}^{(i)}$ by using a 3-layer network. \\
8. \> Predict the target bits of test sample ${\bf t}^{\prime}$ from the trained \\
   \> network. (see Algorithm 5)
\end{tabbing}

To enhance the capability of proposed method, two types of distributions of random points in each sample 
are synthesized. The first type of distribution is based on Gaussian distribution with a randomly 
predefined standard deviation value. The second type of distribution is based on random distribution.
Let $x^{(i)}_{j,k}$ be a coordinate of point ${\bf p}^{(i)}_{j}$ in dimension $k$ of sample ${\bf s}_{i}$. 
A center is also considered as a data point.

\vspace{0.2in}
\noindent
\IEEEpubidadjcol
{\bf Algorithm 2: Synthesizing a pre-learning sample ${\bf s}_{i}$}
\label{alg2}
\begin{tabbing}
1. \= Generate a set of random cluster centers, ${\bf Z} = \{{\bf c}_{1}, \ldots, $ \\
   \> ${\bf c}_{n}\}$; $1 \leq n \leq C$, in terms of coordinates in $d$ dimensions. \\
2. \> Generate a random integer $m$ such that $N \leq m \leq M$. \\
3. \> Generate a random standard deviation value $\sigma$. \\
4. \> Alternatively use steps 4.1 and 4.2 to generate random \\
   \> centers and points.\\
   \> 4.1 \= Generate a set of random points ${\bf P}^{(i)}$ with respect to \\
   \>     \> all centers in ${\bf Z}$ with a random standard deviation $\sigma$ \\
   \>     \> such that $|{\bf P}^{(i)}| = m$. \\
   \> 4.2 \> Generate a set of random points ${\bf P}^{(i)}$ with respect to \\
   \>     \> all centers in ${\bf Z}$ with a random standard deviation $\sigma$ \\
   \>     \> such that $|{\bf P}^{(i)}| = m$. \\
5. \> Let $L_{k} = \min\limits_{\forall {\bf p}_{i}} (x_{i,k})$. and  $U_{k} = \max\limits_{\forall 
{\bf p}_{i}} (x_{i,k})$. \\
\IEEEpubidadjcol
6. \> Normalize $ x^{(i)}_{j,k} = $ $\left\{ \begin{array}{ll}
                                 \frac{x^{(i)}_{j,k} + |L_{k}|}{|U_{k}|} & \mbox{if $L_{k} < 0$} \\
                                 \frac{x^{(i)}_{j,k}}{|U_{k}|} & \mbox{if $L_{k} > 0$}
                                 \end{array} 
                                 \right . $
\end{tabbing}

All coordinates of each sample ${\bf s}_{i}$ are normalized to lie within the same regional space. This 
is beneficial to achieve the high accuracy of leaning process for all samples.
Obviously, the number of points in each sample after Algorithm 2 may not be equal. This will make the 
result of transforming the point distribution of each sample into a vector of equal dimension prior to the
learning process impossible. Thus, Algorithm 3 is used to synthesize some extra points. It is important 
that these extra points must not interfere and disturb the original topology of point distribution in each 
sample.

\vspace{0.2in}
\noindent
{\bf Algorithm 3: Making all sample sizes equal and transforming point distribution into a vector}
\begin{tabbing}
{\bf Input:} \= 1. \= A sample ${\bf s}_{i}$ whose set of points is ${\bf P}^{(i)} = \{{\bf p}^{(i)}_{1},$\\ 
             \>    \> ${\bf p}^{(i)}_{2}, \ldots, {\bf p}^{(i)}_{m}\}$ such that ${\bf p}^{(i)}_{j}  = 
                   [x^{(i)}_{j,1}, x^{(i)}_{j,2}, \ldots, $ \\
             \>    \> $x^{(i)}_{j,d}]^{T}$. \\
             \> 2. \> Maximum number of points, $M$.
\end{tabbing}
\begin{tabbing}
{\bf Output:} \= 1. \= Sample $|{\bf s}_{i}| = M$ with extra points.\\
              \> 2. \> ${\bf v}_{i} = [a^{(i)}_{1}, a^{(i)}_{2}, \ldots, a^{(i)}_{L}]^{T}$, where $L = d 
                    \times |{\bf P}^{(i)}|$.
\end{tabbing}

\begin{tabbing}
1. \= {\bf For} \= $1 \leq j \leq M - |{\bf s}_{i}|$:  \\
2. \>           \> Randomly pick a point ${\bf p}^{(i)}_{j} \in {\bf P}^{(i)}$. \\
3. \>           \> Set ${\bf p}^{(i)}_{i+m} = {\bf p}^{(i)}_{j}$. \\
4. \>           \> ${\bf P}^{(i)} = {\bf P}^{(i)} \cup \{{\bf p}^{(i)}_{i+m}\}$. \\
5. \> Let ${\bf v}^{(i)}_{j}$ be an empty sample vector for sample ${\bf s}_{i}$. \\
6. \> {\bf For} \= $1 \leq k \leq d$: \\
7. \>           \> Sort $x^{(i)}_{j,k}; \: 1 \leq j \leq |{\bf P}^{(i)}|$ in an ascending order. \\
8. \>           \> Append all sorted $x^{(i)}_{j,k}; \: 1 \leq j \leq |{\bf P}^{(i)}|$ to ${\bf v}_{i}$. 
\end{tabbing}

After generating extra points to make the size of sample ${\bf s}_{i}$ equal $M$, all distributed points 
are transformed into a vector of size $d \times M$. Since each point is at a random location in a $d$ 
dimensional space, it is rather difficult to order these points according to their locations. Thus, the 
coordinate in each dimension of each point is considered instead. All coordinates of all points in each 
dimension are collected and sorted in an ascending order. Then all these sorted coordinates in all 
dimensions are concatenated to form a vector in $R^{d \times M}$, as appeared in steps 6-8 in Algorithm 3.

\vspace{0.2in}
\noindent
{\bf Algorithm 4: Self-encoding sample vectors and testing vector}
\begin{tabbing}
{\bf Input:}  \= 1. \= A set of all sample vectors, ${\bf V} = \{{\bf v}_{1} {\bf v}_{2}, \ldots, {\bf v}_{M}\}$ \\ 
              \>    \> such that where $L = d \times |{\bf P}_{i}|$. \\
              \> 2. \> ${\bf v}_{i} = [a^{(i)}_{1}, a^{(i)}_{2}, \ldots, ^{(i)}{L}]^{T}$, where $L = d 
                       \times |{\bf P}^{(i)}|$. \\
              \> 3. \> A test sample ${\bf t}= [b_{1}, b_{2}, \ldots, b_{L}]^{T}$.
\end{tabbing}
\begin{tabbing}
{\bf Output:} \= All encoded sample vectors from the output of \\
              \> hidden layer.
\end{tabbing}
\begin{tabbing}
1. ~\= Create a 3-layer network whose input layer has $L$ neurons, \\
   ~\> hidden layer has 1 neuron, and output layer has $L$ neurons. \\
2. ~\> Let ${\bf V} = {\bf V} \cup \{{\bf t}\}$. \\
3. ~\> Let $o_{j}$ be the output value of output neuron $j$. \\
4. ~\> Compute $e = \sum_{i=1}^{M+1} \sum_{k=1}^{L} (a^{(i)}_{k} - o_{k})^{2}$. \\
5. ~\> Let $\epsilon$ be a error threshold and $T$ be the maximum \\
   ~\> number of epochs. \\
6. ~\> {\bf While} \= $e > \epsilon$ {\bf and} iterations $< T$: \\
7. ~\>             \> Add 1 hidden neuron. \\
8. ~\>             \> Retrain the network. \\
9. ~\>             \> Compute $e = \sum_{i=1}^{M+1} \sum_{k=1}^{L} (a^{(i)}_{k} - o_{k})^{2}$. \\
10.~\> {\bf End} 
\end{tabbing}
Suppose the number of hidden neurons after the encoding process is $h$. Let ${\bf V}^{\prime} = 
\{{\bf v}^{\prime}_{1}, {\bf v}^{\prime}_{2}, \ldots, {\bf v}^{\prime}_{M}, {\bf v}^{\prime}_{M+1}\}$ 
be the set of encoded sample vectors from Algorithm 4. The dimensions of each encoded vector are equal 
to $h$. The last ${\bf v}_{M+1}$ is the encoded test vector. These $M$ encoded vectors are used to train 
a 3-layer neural network. Then the target of encoded test vector ${\bf v}^{\prime}_{M+1}$ is predicted 
by using this trained network. The detail is in Algorithm 5. 

\vspace{0.2in}
\noindent
{\bf Algorithm 5: Predicting number of clusters of testing data}
\begin{tabbing}
{\bf Input:}  \= 1. \= A set of all encoded training vectors, ${\bf V}^{\prime} = $ \\
              \>    \> $\{{\bf v}^{\prime}_{1}, {\bf v}^{\prime}_{2}, \ldots, {\bf v}^{\prime}_{M}\}$ \\
              \> 2. \> A testing vector ${\bf v}^{\prime}_{M+1}$. \\
              \> 3. \> A target ${\bf t}_{i} = [b_{i,1}, b_{i,2}, \ldots, b_{i,C}]^{T}$. $C$ for each 
                       input \\
              \>    \> vector ${\bf v}^{\prime}_{i}$, including both training and test vectors. $C$ \\
              \>    \> is the upper bound of number of clusters.
\end{tabbing}
\begin{tabbing}
{\bf Output:} \= Predicted number of clusters.
\end{tabbing}
\begin{tabbing}
1. ~\= {\bf For} \= each training vector ${\bf v}^{\prime}_{i}$ {\bf do} \\
2. ~\>           \> Set target bit of $b_{i,j} = 0.5$ corresponding to the \\
    \>           \> number of clusters $j$ generated by Algorithm 2. \\
3. ~\> Construct a 3-layer feed-forward network with $h$ inputs \\
    \> and $C$ outputs. \\
4. ~\> Training the network until the target error is less than a \\
    \> predefined value. \\
5. ~\> Let $b_{M+1,j}$, $1 \leq j \leq C$ be the predicted output of \\
    \> bit $j$ of test vector predicted by the trained network.\\
6. ~\> Select bit $j$ of predicted output such that value of bit j \\
    \> is equal to $\max\limits_{1 \leq j \leq C} (b_{M+1,j})$. \\
7. ~\> Return $j$ as the number of predict clusters.
\end{tabbing}

\subsection{Examples}

\begin{figure*}[t]
\centering
\includegraphics[scale=0.55]{overview-example.eps}
\vspace{-1.6in}
\caption{An example of how the proposed algorithms work for the pre-training process. There are two 
samples. The first sample has 7 data points and the second sample has 9 data points.}
\label{overview-example}
\end{figure*}
\newtheorem{thm}{Theorem}
\newtheorem{lemma}{Lemma}

\section{Theoretical Analysis of Upper Bound of Synthetic Patterns}

Transforming a problem of clustering into a problem of classification to identify the actual
number of clusters may not achieve the absolute accuracy due to the insufficient number of training 
patterns. The interesting problem is how many pattern distributions of fixed number of clusters in $d$ 
dimensional space must be synthesized. Actually, this problem is the same as the problem of sphere 
packing \cite{Cohen-Elkies}\cite{Gaber}\cite{Goldberg}\cite{Torquato-Stillinger}. The problem of sphere 
packing can be briefly addressed as follows. Given a space in a $d$ dimensional Euclidean space, how many 
spheres of different sizes can be packed so that the spheres touch others at their boundaries. In a high 
dimensional space, this problem is very complicated. Instead of indicating the number spheres, the upper 
bound of density defined as the ratio of container volume and total volume of all packed spheres is 
measured. However, estimating the density is not simple in case of hyper-spheres. Gabor Fejes Toth 
\cite{Gaber} summarized the upper bound of packing density in a lattice structure
($\delta_L(d)$) found by several researchers for 2-8, and 24 dimensions. For example, $\delta_L(2) = 
\pi/\sqrt{12}$, $\delta_L(3) = \pi/\sqrt{18}$, $\delta_L(4) = \pi^{2}/16$, and $\delta_L(5) = \pi^{2}/
(15\sqrt{2})$. All circles and spheres have one unit radius.

In this paper, only 2, 3, and 4 dimensions are considered. The constraints imposed on the derivative 
upper bound of pre-learned patterns are the followings.
\begin{enumerate}
\item Each cluster has the same volume size.
\item The shape of each cluster is in a form of a circle (in 2 dimensions) and a sphere (in 3 dimensions).
\item The radius of each hyper-sphere is less than the width of a container.
\item The shape of container is in a form of a square (in 2 dimensions) and a cube (in 3 dimensions).
\end{enumerate}
Figure \ref{packing-density} illustrates the concept of deriving the upper bound of pre-learned patterns 
in a 2-dimensional space. Each cluster is represented by a unit circle.  
The concept consists of two phases. In the first phase, the upper bound of packing density is considered 
and the number of possible circles is computed from the upper bound of packing density. In the second 
phase, the number of pre-learned patterns for each number of clusters is derived. 

In 2 dimensions, let $V$ be the area of a container, $N$ be the number of circles, and $r$ be the radius.
Based on the theoretical upper bound in \cite{Gaber}, the packing density is $\pi/\sqrt{12}$.

\begin{lemma}
In 2 dimensions, the maximum number of packed circles is $N = \frac{0.9069}{\pi r^{2}}$.
\end{lemma} 
\begin{proof}
The packing density can be computed from the ratio of $\frac{N \pi r^{2}}{V}$. The area of container is one
unit. Thus, we have
\begin{eqnarray}
\frac{N \pi r^{2}}{V} & = & \frac{\pi}{\sqrt{12}} \\
N & = & \frac{\pi}{\sqrt{12}} \frac{V}{\pi r^{2}} \\
 & = & \left \lceil \frac{0.9069V}{\pi r^{2}} \right \rceil
\end{eqnarray}
\end{proof}
Since the area of $V$ is one unit, the value of radius $r$ must be less than 1. The number of pre-learned 
patterns must cover the number of possible clusters of a queried data set. Let $C$ be the maximum number 
of clusters. There are two scenarios concerning the value of $C$. For classification problem, the value 
of $C$ is unknown from a data set. In this circumstance, for safety, the value of $C$ can be set to the 
number of classes plus $a$, where $a \geq 1$. For clustering problem, we assume that the shape of clusters 
are governed by some known distribution such as Gaussian distribution. In this situation, defining the 
value of $C$ is as difficult as defining the number of clusters prior to the clustering process. One 
possible value of $C$ is to estimate the value from $N$ in Lemma 1.

\begin{figure}[!t]
\centering
\includegraphics[width=3.0in]{packing-density.eps}
\caption{Concept of lower bound analysis of pre-learned patterns based on the imposed constraints. The possible distribution of clusters are shown by colored circles and transparent circles. The colored
circles denote the locations of clusters and the transparent circles denote that empty locations.}
\label{packing-density}
\end{figure}

\begin{thm}
In 2 dimensions, the upper bound of pre-learned patterns $L$ in a 2-dimensional space is 
\begin{equation}
L = \sum\limits_{c=1}^{C} \left ( \begin{array}{c} 
                                   N \\
                                   c
                                  \end{array} \right )
\end{equation}
\noindent
for $N = \left \lceil \frac{0.9069V}{\pi r^{2}} \right \rceil$.
\end{thm}
\begin{proof}
The reason is from Figure \ref{packing-density}. Each circle denotes the possible location of each cluster 
without any overlap. Hence, the number of possible locations to distribute $c$, $1 \leq c \leq C$ clusters 
in the container is $\left ( \begin{array}{c} N \\ c \end{array} \right )$.
\end{proof} 
In 3 dimensions, the same argument and derivation as that in the 2 dimensional space can be followed.
The theoretical upper bound of packing density is from Kepler conjecture \cite{Hales1} which stated that 
\begin{quote}
{\em No packing of congruent balls in Euclidean three space has density greater than that of the face-
centered cubic packing}
\end{quote}
The upper bound of packing density is $\pi/\sqrt{18} \approx 0.74$. This conjecture was later proved by 
Thomas C. Hales \cite{Hales2}. Therefore, we have the following upper bound of pre-learned patterns for 3 
dimensions.

\begin{lemma}
In 3 dimensions, the maximum number of packed circles is $N = \left \lceil \frac{0.74V}{1.33\pi 
r^{3}} \right \rceil $.
\label{3d-lemma}
\end{lemma}
 
\begin{thm}
In 3 dimensions, the upper bound of pre-learned patterns $L$ in a 3-dimensional space is 
\begin{equation}
L = \sum\limits_{c=1}^{C} \left ( \begin{array}{c} 
                                   N \\
                                   c
                                  \end{array} \right )
\end{equation}
\end{thm}
\noindent
for $N = \left \lceil \frac{0.74V}{1.33\pi r^{3}} \right \rceil$.

\begin{lemma}
In 4 dimensions, the maximum number of packed circles is $N = \left \lceil \frac{0.6169\times2V}{\pi^{2} 
r^{4}} \right \rceil $.
\end{lemma}
 
\begin{thm}
In 4 dimensions, the upper bound of pre-learned patterns $L$ in a 3-dimensional space is 
\begin{equation}
L = \sum\limits_{c=1}^{C} \left ( \begin{array}{c} 
                                   N \\
                                   c
                                  \end{array} \right )
\end{equation}
\end{thm}
\noindent
for $N = \left \lceil \frac{0.6169\times2V}{\pi^{2} r^{4}} \right \rceil$.
%--------------------- not correct -------------------------------
\begin{table}[t]
\caption{Lower bound of pre-learned patterns in 2-4 dimensions for each radius and density \cite{Gaber}. 
The volume is a hyper-square of size $10^{d}$, where $d$ is the number of dimensions.}
\begin{center}
\begin{tabular}{|c|c|c|c|c|c|c|c|}\hline
\multirow{2}{*}{density \cite{Gaber}} & $r =$ & \multicolumn{5}{c|}{\# pre-learned patterns for \# clusters} 
 & \multirow{2}{*}{total} \\ \cline{3-7}
 & $0.75d$ & \#1 & \#2 & \#3 & \#4 & \#5 & \\ \hline
0.9069 (d=2) & 1.5 & 13 & 78 & 286 & 715 & 1287 & 2379 \\ \hline
0.7405 (d=3) & 2.25 & 16 & 120 & 560 & 1820 & 4368 & 6884 \\ \hline
0.6169 (d=4) & 3.0 & 16 & 120 & 560 & 1820 & 4368 & 6884 \\ \hline
\end{tabular}
\end{center}
\end{table}
%-----------------------------------------------------------

\section{Experiments and Results}

The accuracy of estimation was measured from two experiments. The first experiment is based on the 
synthetic patterns in a 2 dimensional space. The second experiment is based on iris data set whose number 
of clusters is equal to 3. The detail of both experiments are in the following sections.

\subsection{Experiments on Synthetic Data} 

The package {\em scikit-learn} version 1.3 was used to generate synthetic data sets with random
standard deviations and random locations of cluster centers in the experiment. 
The number of clusters in each learning pattern ranges from 1 cluster to 5 clusters. There are 700 
synthetic patterns for each number of clusters which makes the total number of patterns equal to $5 \times 
700 = 3500$. To evaluate the accuracy, the experiment was conducted by using 700 folds of cross validation.
For each fold, one randomly selected patterns from 3500 patterns is for testing and the rest is
for training the network. Some of the training patterns are shown in Figure \ref{sythetic-cluster-data} 
and the testing patterns are shown in Figure \ref{test-cluster-data}. There are 685 folds giving the 
correct number of clusters, which is equal to 98\%. Figure \ref{wrong-test-cluster-data} shows the 
clustering patterns of wrong prediction. Since the testing patterns were randomly selected, it is
possible that some patterns were repeatedly selected.

\begin{figure*}[t]
\begin{center}
\begin{tabular}{c@{\hskip 0.1in}c@{\hskip 0.1in}c@{\hskip 0.1in}c@{\hskip 0.1in}c@{\hskip 0.1in}}
\includegraphics[scale=0.2]{train-cluster1-dataset1.eps} & 
\includegraphics[scale=0.2]{train-cluster1-dataset2.eps} &
\includegraphics[scale=0.2]{train-cluster1-dataset3.eps} &
\includegraphics[scale=0.2]{train-cluster1-dataset4.eps} &
\includegraphics[scale=0.2]{train-cluster1-dataset5.eps} \\
 & & (a) 1 cluster. & & \\
\includegraphics[scale=0.2]{train-cluster2-dataset1.eps} & 
\includegraphics[scale=0.2]{train-cluster2-dataset2.eps} &
\includegraphics[scale=0.2]{train-cluster2-dataset3.eps} &
\includegraphics[scale=0.2]{train-cluster2-dataset4.eps} &
\includegraphics[scale=0.2]{train-cluster2-dataset5.eps} \\
 & & (b) 2 clusters. & & \\
\includegraphics[scale=0.2]{train-cluster3-dataset1.eps} & 
\includegraphics[scale=0.2]{train-cluster3-dataset2.eps} &
\includegraphics[scale=0.2]{train-cluster3-dataset3.eps} &
\includegraphics[scale=0.2]{train-cluster3-dataset4.eps} &
\includegraphics[scale=0.2]{train-cluster3-dataset5.eps} \\
 & & (c) 3 clusters. & & \\
\includegraphics[scale=0.2]{train-cluster4-dataset1.eps} & 
\includegraphics[scale=0.2]{train-cluster4-dataset2.eps} &
\includegraphics[scale=0.2]{train-cluster4-dataset3.eps} &
\includegraphics[scale=0.2]{train-cluster4-dataset4.eps} &
\includegraphics[scale=0.2]{train-cluster4-dataset5.eps} \\
 & & (d) 4 clusters. & & \\
\includegraphics[scale=0.2]{train-cluster5-dataset1.eps} & 
\includegraphics[scale=0.2]{train-cluster5-dataset2.eps} &
\includegraphics[scale=0.2]{train-cluster5-dataset3.eps} &
\includegraphics[scale=0.2]{train-cluster5-dataset4.eps} &
\includegraphics[scale=0.2]{train-cluster5-dataset5.eps} \\
 & & (e) 5 clusters. & & \\
\end{tabular}
\end{center}
\caption{Some examples of synthetic data for learning as past knowledge. (a) One cluster. (b) Two clusters. 
(c) Three clusters. (d) Four clusters. (e) Five clusters.}
\label{sythetic-cluster-data}
\end{figure*}

%-------------------------------------------------
\begin{figure*}[t]
\begin{center}
\begin{tabular}{c@{\hskip 0.1in}c@{\hskip 0.1in}c@{\hskip 0.1in}c@{\hskip 0.1in}c@{\hskip 0.1in}}
\includegraphics[scale=0.2]{test-train-cluster1-dataset132.eps} & 
\includegraphics[scale=0.2]{test-train-cluster1-dataset134.eps} &
\includegraphics[scale=0.2]{test-train-cluster1-dataset170.eps} &
\includegraphics[scale=0.2]{test-train-cluster1-dataset312.eps} &
\includegraphics[scale=0.2]{test-train-cluster1-dataset670.eps} \\
 & & (a) 1 cluster. & & \\
\includegraphics[scale=0.2]{test-train-cluster2-dataset124.eps} & 
\includegraphics[scale=0.2]{test-train-cluster2-dataset364.eps} &
\includegraphics[scale=0.2]{test-train-cluster2-dataset425.eps} &
\includegraphics[scale=0.2]{test-train-cluster2-dataset586.eps} &
\includegraphics[scale=0.2]{test-train-cluster2-dataset652.eps} \\
 & & (b) 2 clusters. & & \\
\includegraphics[scale=0.2]{test-train-cluster3-dataset98.eps} & 
\includegraphics[scale=0.2]{test-train-cluster3-dataset186.eps} &
\includegraphics[scale=0.2]{test-train-cluster3-dataset249.eps} &
\includegraphics[scale=0.2]{test-train-cluster3-dataset527.eps} &
\includegraphics[scale=0.2]{test-train-cluster3-dataset540.eps} \\
 & & (c) 3 clusters. & & \\
\includegraphics[scale=0.2]{test-train-cluster4-dataset148.eps} & 
\includegraphics[scale=0.2]{test-train-cluster4-dataset305.eps} &
\includegraphics[scale=0.2]{test-train-cluster4-dataset337.eps} &
\includegraphics[scale=0.2]{test-train-cluster4-dataset410.eps} &
\includegraphics[scale=0.2]{test-train-cluster4-dataset607.eps} \\
 & & (d) 4 clusters. & & \\
\includegraphics[scale=0.2]{test-train-cluster5-dataset140.eps} & 
\includegraphics[scale=0.2]{test-train-cluster5-dataset214.eps} &
\includegraphics[scale=0.2]{test-train-cluster5-dataset282.eps} &
\includegraphics[scale=0.2]{test-train-cluster5-dataset339.eps} &
\includegraphics[scale=0.2]{test-train-cluster5-dataset596.eps} \\
 & & (e) 5 clusters. & & \\
\end{tabular}
\end{center}
\caption{Some examples of test data whose number of clusters is correctly predicted. (a) One cluster. (b) 
Two clusters. (c) Three clusters. (d) Four clusters. (e) Five clusters.}
\label{test-cluster-data}
\end{figure*}

%----------------------------------------------------------------------------
\begin{figure*}[t]
\begin{center}
\begin{tabular}{c@{\hskip 0.1in}c@{\hskip 0.1in}c@{\hskip 0.1in}}
\includegraphics[scale=0.2]{wrong-test-train-cluster1-dataset195.eps} & 
\includegraphics[scale=0.2]{wrong-test-train-cluster3-dataset572.eps} &
\includegraphics[scale=0.2]{wrong-test-train-cluster4-dataset602.eps} \\
(a) 1 cluster but predicted 2 clusters. & (b) 3 clusters but predicted 2 clusters. & (c) 4 clusters
but predicted 2 clusters.
\end{tabular}
\end{center}
\caption{Some examples of test data whose number of clusters is incorrectly predicted. (a) 1 cluster. (b)
2 clusters. (c) 4 clusters.}
\label{wrong-test-cluster-data}
\end{figure*}

\subsection{Experiment on Real Data Set}

The proposed method was applied to predict the number of clusters in iris data set. This data set is in 4 
dimensions which makes it difficult to visualize the correct number of clusters. However, the number of 
clusters can be presumed to be equal to 3, which is the number of classes. In the experiment, the number
of possible clusters was set from 1 to 5. For each number of clusters, 700 pre-training patterns in
4 dimensions were synthesized. The proposed algorithm predicted 2 clusters in iris data set which is less
than the number of classes. To empirically verify whether this result is correct or not, a set of scatter
plots in 3 dimensions were visualized as illustrated in Figure \ref{iris-data}.
It can be seen that that data of Iris-virginica overlap data of Iris-versicolor classes. This is the 
reason why the number of predicted clusters is equal to 2.

\begin{figure}[ht!]
\begin{center}
\begin{tabular}{cc}
\includegraphics[scale=0.27]{iris-012.eps} & \includegraphics[scale=0.27]{iris-013.eps} \\
\includegraphics[scale=0.27]{iris-023.eps} & \includegraphics[scale=0.27]{iris-123.eps}
\end{tabular}
\end{center}
\caption{Four combinations of three features (dimensions) of Iris data distributions. Blue points denote 
Iris-setosa. Orange points denote Iris-versicolor. Green points denote Iris-virginica. Some of 
Iris-verginica data points overlap Iris-versicolor data points.}
\label{iris-data}
\end{figure}

\section{Discussion}
This new approach does not require any statistical analysis to predict the number of clusters (k-value)
prior to the actual process of identifying the location of each cluster as well as the members of cluster.
Although the preliminary experiments significantly imply the merit of this proposed methods, the accuracy
depends on two factors which are the set of pre-training data and patterns of data distribution.
The first factor which concerns the upper bound of number of training data can be theoretically estimated
in the forms of sphere packing in a high dimensional space. However, the current theoretical bound of
high dimensional packing is a very difficult problem which makes the estimation of maximum number of
training data become an open problem. The second factor can be alleviated by synthesizing the training
data with all available statistical distributions and also with some artificial distribution such as 
nested spiral shapes. This artificial distribution, in fact, does not exist in nature. It is just used 
for benchmarking the performance of proposed methods. Therefore, in practice, the proposed method in 
this paper is suitable for coping with real data, whose distribution can be imitated by some well-defined 
statistical distributions. But it is still possible to apply the proposed method to the artificial data
as well. 

New synthetic data distribution patterns can be continuously added and learned by the pre-trained neural 
network without involving any previously trained data. One possible approach to this type of training 
is by applying the concept of {\em discard-after-learn} \cite{Prem}\cite{Mongkhon} whose the training time 
complexity is $O(P^{2})$ where $P$ is the number of training patterns. 

For our approach, the time 
complexity of the training process can be analyzed as follows. The data preparation before encoding step 
is $O(NP)$ where $N = \max\limits_{i=1,P}(n_{i})$; $P$ is the number of training patterns and $n_{i}$ is the 
number of data points in each training pattern $i$. The encoding input patterns and training of encoded 
pattern steps were trained by a feedforward network implemented in Tensorflow and Keras environments.
Unfortunately, the time complexity of these two steps is rather complex to analyze.

\section{Conclusion}

A new approach to identify the number of clusters (k-value) without applying any statistical distance 
analysis is proposed. This method is based on the finding in cognitive science regarding the study of how 
the brain uses past knowledge and experiences to solve problems. The problem of identifying number of
clusters is transformed into a problem of classifying data distribution patterns according to the
actual number of clusters.  The preliminary study and experiments signify this proposed method.
The accuracy of prediction depends on the number of pre-trained cluster patterns with various 
distributions. Further study should focus on how to estimate the upper bound of number of pre-trained 
patterns and how to synthesize various distributions with some degree of overlap involved in a high 
dimensional space.

\begin{thebibliography}{1}
\bibliographystyle{IEEEtran}

\bibitem{Hong}
Hong Yu, Zhanguo Liu, Guoyin Wang, ``An automatic method to determine the number of clusters using
decision-theoretic rough set'', International Journal of Approximate Reasoning 55 (2014) 101–115.

\bibitem{Calmon}
Wilson Calmon and Mariana Albi, ``Estimation the number of clusters in a ranking data context'', Information Science
546(2021) 977-995.

\bibitem{Abdul} 
Md Abdul Masud, Joshua Zhexue Huang, Chenghao Wei, Jikui Wang, Imran Khan, Ming Zhong, ``I-nice: A new approach 
for identifying the number of clusters and initial cluster centres'', Information Science 466 (2018) 129–151.

\bibitem{Liang}
Jiye Liang, Xingwang Zhao, Deyu Li, Fuyuan Cao, Chuangyin Dang, ``Determining the number of clusters using 
information entropy for mixed data'', Pattern Recognition 45 (2012) 2251–2265.

\bibitem{Gupta}
Avisek Gupta, Shounak Datta, Swagatam Das, ``Fast automatic estimation of the number of clusters from the
minimum inter-center distance for k-means clustering'', Pattern Recognition Letters 116 (2018) 72–79.

\bibitem{Gialampoukidis}
Ilias Gialampoukidis, Stefanos Vrochidis, Ioannis Kompatsiaris, Ioannis Antoniou, ``Probabilistic 
density-based estimation of the number of clusters using the DBSCAN-martingale process'', Pattern 
Recognition Letters 123 (2019) 23–30.

\bibitem{Chaea}
Seong S. Chaea, Janice L. DuBienb, William D. Warde, ``A method of predicting the number of clusters 
using Rand’s statistic'', Computational Statistics \& Data Analysis 50 (2006) 3531–3546.

\bibitem{Bhowmik}
Mrinal Kanti Bhowmik, Tathagata Debnath, Debotosh Bhattacharjee, Paramartha Dutta, ''EF-Index: Determining 
number of cluster (K) to estimate number of segments (S) in an image'', Image and Vision Computing 
88(2019) 29-40.

\bibitem{Specht}
E. Specht, ``Hugh Densitu packings of equal circles in rectangles with variable aspect ration'', 
Computers \& Operations Research 40(2013) 58-69.

\bibitem{Lippiello}
E. Lippiello, S. Baccari, P. Bountzis, ``Determining the number of clusters, before finding clusters,
from the susceptibility of the similarity matrix'', Physica A 616 (2023) 128592.

\bibitem{Wang}
Pingxin Wang, Yiyu Yao, ``CE3: A three-way clustering method based on mathematical morphology'',
Knowledge-Based Systems 155 (2018) 54–65.

\bibitem{Ramazan}
Ramazan Ünlü, Petros Xanthopoulos, ``Estimating the number of clusters in a dataset via consensus 
clustering'', Expert Systems With Applications 125 (2019) 33–39.

\bibitem{Fahim}
Ahmed Fahim, ``K and starting means for k-means algorithm'', Journal of Computational Science 55 (2021) 101445.

\bibitem{Luo}
Wei Luo, ``Determine the number of clusters by data augmentation'', Electronic Journal of Statistics
Vol. 16 (2022) 3910–3936.

\bibitem{He}
Zhenzhen He, Zongpu Jia,and  Xiaohong Zhang, ``A Fast Method for Estimating the Number of Clusters
Based on Score and the Minimum Distance of the Center Point'', Information 2020, 11, 16.

\bibitem{Toan}
Toan Nguyen Mau, Yasushi Inoguchi, and Van-Num Huynh, ``A Novel Cluster Prediction Approach Based
on Locality-Sensitive Hashing for Fuzzy Clustering of Categorical Data, '', IEEE Access, vol. 10, 2022, 
pp. 34196-34206.

\bibitem{Xiaolong}
Xiaolong Wang, Yiping Jiao, Shumin Fei, ``Estimation of Clusters Number and Initial Centers of K-means
Algorithm Using Watershed Method'', 2015 14th International Symposium on Distributed Computing and Applications 
for Business Engineering and Science, pp. 505-508.

\bibitem{Rahman}
Md Moshiur Rahman, Abdul Masudy, Badhan Mazumderz, ``Estimation of the Number of Clusters based on
Simplical Depth'', 2020 2nd International Conference on Sustainable Technologies for
Industry 4.0 (STI).

\bibitem{Dong}
Dong Sik Kim, ``Finding the Number of Clusters Using a Small Training Sequence'', IEEE Access, Vol. 11, 2023,
pp. 25932-25940.
  
\bibitem{Cohen-Elkies}
H. Cohen and N. Elkies, ``New upper bounds on sphere packings I'', \textit{Annals of Mathematics}, 157 (2003), 
689–714.

\bibitem{Gaber}
Gabor Fejes Toth, ``Packing and Covering'', \textit{Handbook of Discrete and Computational Geometry},
J.E. Goodman, J. O'Rourke, and C. D. Tóth (editors), 3rd edition, CRC Press, Boca Raton, FL, 2017

\bibitem{Goldberg}
M. Goldberg, ``The packing of equal circles in a square'', \textit{Mathematics Magazine}, Jan., 1970, 
Vol. 43, No. 1 (Jan., 1970), pp. 24-30.

\bibitem{Torquato-Stillinger}
S. Torquato and F.H. Stillinger, ``New conjectural lower bounds on the optimal density of sphere 
packings'', Experimental Mathematics, Vol. 15 (2006), No. 3. pp. 307-331.
 
\bibitem{Karoly}
K. Bezdek, ``Circle packings into convex domains of the Euclidean and hyperbolic plane and the sphere'', 
Geometriae Dedicata 21 (1968) pp. 249-255.

\bibitem{Hales1}
Thomas C. Hales, ``A proof of the Kepler conjecture'', Annals of Mathematics, 162 (2005), 1065–1185.

\bibitem{Hales2}
Thoms C. Hales, ``The sphere packing problem'', Journal of Computational and Applied Mathematics 44(1992) 41-76.

\bibitem{Prem}
Prem Junsawang, Suphakant Phimoltares, Chidchanok Lursinsap, ``A fast learning method for streaming and 
randomly ordered multi-class data chunks by using one-pass-throw-away class-wise learning concept'',
Expert Systems With Applications 63 (2016) 249–266.

\bibitem{Mongkhon}
Mongkhon Thakong, Suphakant Phimoltares, Saichon Jaiyen, and Chidchanok Lursinsap, ``Fast Learning and 
Testing for Imbalanced Multi-Class Changes in Streaming Data by Dynamic Multi-Stratum Network'',
IEEE Access, Volume 5, 2017, pp. 10633-10648.

\bibitem{Dagmar}
Dagmar Divijak, ``Predicting: Using Past Experience to Guide Future Action'', Chapter 8, Frequency in 
Language, Cambridge University Press, September 2019.

\bibitem{Natalie}
Natalie Biderman, Akram Bakkour,and Daphna Shohamy, ``What Are Memories For? The Hippocampus
Bridges Past Experience with Future Decisions'', Trends in Cognitive Sciences, July 2020, Vol. 24, No. 7,
pp. 542-556.

\bibitem{Hawkins}
J. Hawkins, and S. Blakeslee, On Intelligence, 2004 (Henry Holt: New York).

\bibitem{Bar}
M. Bar, Predictions in the Brain: Using Our Past to Generate a Future, 2011 (Oxford
University Press: New York, Oxford)
%---------------------------------


\end{thebibliography}


\newpage

%\section{Biography Section}
%If you have an EPS/PDF photo (graphicx package needed), extra braces are
% needed around the contents of the optional argument to biography to prevent
% the LaTeX parser from getting confused when it sees the complicated
% $\backslash${\tt{includegraphics}} command within an optional argument. (You can create
% your own custom macro containing the $\backslash${\tt{includegraphics}} command to make things
% simpler here.)
 
%\vspace{11pt}

%\bf{If you include a photo:}\vspace{-33pt}
%\begin{IEEEbiography}[{\includegraphics[width=1in,height=1.25in,clip,keepaspectratio]{fig1}}]{Michael Shell}
%Use $\backslash${\tt{begin\{IEEEbiography\}}} and then for the 1st argument use $\backslash${\tt{includegraphics}} to declare and link the author photo.
%Use the author name as the 3rd argument followed by the biography text.
%\end{IEEEbiography}

\vspace{11pt}

%\bf{If you will not include a photo:}\vspace{-33pt}
%\begin{IEEEbiographynophoto}{John Doe}
%Use $\backslash${\tt{begin\{IEEEbiographynophoto\}}} and the author name as the argument followed by the biography text.
%\end{IEEEbiographynophoto}




%\vfill

\end{document}


